\section{Discussion and Conclusion}
\label{sec:discussion}

\subsection{Theoretical Contributions}

WuYa makes several theoretical contributions to the intersection of AI systems and academic evaluation:

\begin{enumerate}
    \item \textbf{Principles-to-Architecture Mapping}: We demonstrate a systematic methodology for transforming abstract philosophical principles into operational agent architectures. This provides a template for other domains where theoretical frameworks exist but have not been computationally operationalized.

    \item \textbf{Retrieval-Evaluation Coupling}: We introduce the principle that retrieval should not be merely supplementary but should serve as the core triggering mechanism for evaluation. This reframes \rag{} from knowledge lookup to reasoning catalyst.

    \item \textbf{Self-Improvement for Evaluation}: We extend self-improving agent frameworks from verifiable tasks (code generation, games) to subjective evaluation domains. This opens new research directions at the intersection of agent learning and scholarly communication.
\end{enumerate}

\subsection{Engineering Contributions}

WuYa provides reusable architectural patterns:

\begin{enumerate}
    \item \textbf{Two-Phase Routing}: The separation of normative gatekeeping from substantive evaluation ensures that ethical considerations are addressed before quality assessment. This pattern is applicable to other domains requiring both compliance checking and quality evaluation.

    \item \textbf{Hybrid Knowledge Strategy}: The combination of internalized principles with retrieved evidence balances consistency and flexibility. This architecture can be adapted for other expert domains where both theoretical grounding and contextual knowledge are needed.

    \item \textbf{Frontier Surface Library}: The concept of learned evaluation preferences stored as structured patterns provides a new form of organizational memory for AI systems.
\end{enumerate}

\subsection{Limitations and Future Work}

WuYa has several limitations:

\begin{enumerate}
    \item \textbf{LLM Dependency}: The system depends on large language model APIs, introducing costs, latency, and potential service unavailability. Future work should explore local model deployment with comparable performance.

    \item \textbf{Cold Start for New Journals}: The \dea{} analysis requires historical paper data; new journals without publication history cannot benefit from Path B. The frontier discovery mechanism can help but requires accumulation time.

    \item \textbf{Language Support}: Currently limited to English papers. Extension to multilingual evaluation would require cross-lingual retrieval and evaluation frameworks.

    \item \textbf{Bias Concerns}: Discipline prior distributions reflect existing publication patterns, which may encode historical biases. Careful auditing is needed to ensure WuYa does not perpetuate systemic inequities.
\end{enumerate}

Future directions include:
\begin{itemize}
    \item Integration with manuscript submission systems for seamless workflow.
    \item Personalized recommendation based on author career stage and goals.
    \item Interactive review mode where authors can query WuYa iteratively about improvement strategies.
    \item Cross-disciplinary impact analysis: tracking how research influences other fields.
\end{itemize}

\subsection{Ethical Considerations}

WuYa is designed as an \textit{augmentation} tool for human decision-making, not a replacement. Key ethical principles:

\begin{itemize}
    \item \textbf{Transparency}: Every evaluation includes explanations grounded in canonical texts, enabling human oversight.
    \item \textbf{Fairness}: The self-improvement mechanism should be monitored for bias; regular audits of frontier surfaces are recommended.
    \item \textbf{Privacy}: Paper content and evaluation data should be handled according to data protection regulations.
    \item \textbf{Human-in-the-Loop}: Final decisions on acceptance and journal selection remain with human authors and editors.
\end{itemize}

\subsection{Conclusion}

This paper presented WuYa, a theory-driven multi-agent system that transforms classical philosophical principles into an operational architecture for academic paper evaluation. Through the integration of retrieval-augmented reasoning, two-phase routing, discipline prior calibration, \dea{} efficiency analysis, and self-improving frontier discovery, WuYa achieves significant improvements over existing approaches in both submission recommendation and journal review quality.

WuYa represents a step toward AI systems that are not merely technically sophisticated but also theoretically grounded and continuously learning. By demonstrating that retrieval and self-improvement can be applied to the subjective domain of academic evaluation, we open new research directions at the intersection of AI, philosophy of science, and scholarly communication.

\section*{Acknowledgments}

We thank the anonymous reviewers for their valuable feedback. This work was supported by [funding agency].

\section*{References}
